\chapter{Sources, and how much of each was read}

This chapter names every source the research paper rests on and states how much of it is read, so
that no claim here traces to a document nobody opened.

A surname in a sentence buys a reader nothing; an author, a title and an identifier they can follow
buy them the source. Fields a document does not carry stay empty, with a note naming which. The
extent of reading belongs beside each entry. Five sources are opened to the first page, two are read
in part with the part named, one is read in full, and one cannot be read as text. Without that, a
bibliography degrades into a list of things that exist.

\section{Read in full}

\textbf{Fefferman}~\cite{src:Fefferman}, \emph{Existence and Smoothness of the Navier--Stokes
Equation}. Six pages, errata included, read cover to cover; the Navier--Stokes chapter turns on it.
No year appears anywhere on it, and the paper examined in that chapter cites it without one too; the
entry carries none. Its bytes are identical to the copy claymath.org serves, accessed 11
September 2026.

\section{Read in part, with the part named}

\textbf{Finite time blowup for Navier--Stokes}~\cite{src:OpenAI-2026-Navier-Stokes}. 166 pages.
Abstract, introduction, Theorem 1.1, the bridging argument in section 10.4, Corollary 10.6. Sections
3 through 9 and all three appendices carry the mathematics and are not read.

\textbf{Finite time blowup for the Euler equation}~\cite{src:OpenAI-2026-Euler}. 57 pages.
Abstract, introduction, Theorem 1.1. The amplification mechanism is in the unread remainder, and the
limit test in the Navier--Stokes chapter wants that mechanism.

\textbf{Formal Conjectures}~\cite{src:Formal-Conjectures}, the Lean formalization of the
Navier--Stokes problem statement. Read as source text, in the adapted copy the Lean
repository~\cite{src:OpenAI-2026-Lean} carries, at the commit that file names. Four files out of
2{,}669, none of them built. Everything this research paper finds about the formalized statement
rests on that one file.

\section{Read at the first page only}

Five official Clay problem descriptions, for the five problems the index lists as unsolved. Each is
opened, identified from its own first page instead of its filename, and closed. All five are
byte-identical to the copies claymath.org serves, accessed 11 September 2026.

\textbf{Cook}~\cite{src:Cook}, \emph{The P versus NP Problem}, 12 pages.

\textbf{Deligne}~\cite{src:Deligne}, \emph{The Hodge Conjecture}, 5 pages.

\textbf{Bombieri}~\cite{src:Bombieri}, \emph{Problems of the Millennium: the Riemann Hypothesis}, 11
pages. Older tooling than the other four.

\textbf{Jaffe and Witten}~\cite{src:Jaffe-and-Witten}, \emph{Quantum Yang--Mills Theory}, 14 pages.

\textbf{Wiles}~\cite{src:Wiles}, \emph{The Birch and Swinnerton-Dyer Conjecture}, 5 pages.

One page of Wiles is the largest shortfall here, because Birch and Swinnerton-Dyer is the
problem the toolkit chapter recommends. That recommendation stands on an exact numerical identity,
and the identity is not taken from this document. Whoever acts on the recommendation reads this
first.

\textbf{Wilkins}~\cite{src:Wilkins-1998}, Riemann's memoir translated by David R. Wilkins, \emph{On
the Number of Prime Numbers less than a Given Quantity}, 10 pages. The year is the translation's,
printed there as a preliminary version of December 1998; the memoir's date, November 1859, is
printed beside it. English, legible, and the practical way into the 1859 content.

\section{No text to read}

\textbf{Riemann 1859}~\cite{src:Riemann-1859}, handwritten working manuscript, \emph{Nachlass} Cod.
Ms. Riemann 3, \emph{Rechenseiten}. Four pages with zero extractable characters, counted page by
page, rendered and read as images, the only way they can be read. The year on the entry is the
manuscript's; the scan dates from 2015. Its filename names the published memoir and it holds a
draft, and the corpus chapter states what a citation built from that filename would do.

\section{Not a document}

\textbf{Clay Mathematics Institute}~\cite{src:Clay-Mathematics-Institute}. Two facts in this
research paper describe what that website says on 11 September 2026. The Millennium Problems index
lists five unsolved, one solved, and Navier--Stokes in neither list. The Navier--Stokes page itself
carries the label \emph{Active}. The two do not agree. Both are recorded as statements about how a
website sorts its own contents on one day at one address, and neither is used as a statement about a
problem's status.

\section{Cited for one fact each}

The computations of Riemann zeros~\cite{src:Gourdon-2004,src:Odlyzko-2001,src:Platt-and-Trudgian-2021},
the Birch and Swinnerton-Dyer computation~\cite{src:Birch-and-Swinnerton-Dyer-1965}, the curve
tables~\cite{src:Cremona-1997}, Dokchitser's algorithm~\cite{src:Dokchitser-2004} and the press
report on the prize~\cite{src:RITS-2026} are each cited for the single fact the toolkit or
Navier--Stokes chapter takes from them, as confirmed against the publisher or an index, and not for
their content beyond it.

\section{Named in these sources and not read here}

Reference lists carry names, and this research paper repeats a handful at second hand. Leray on global
weak solutions. Caffarelli, Kohn and Nirenberg on partial regularity. The Beale--Kato--Majda
continuation criterion. Buckmaster and Vicol on non-uniqueness, then Albritton, Bru\'e and Colombo on
the same. Elgindi on axisymmetric blowup.

None is read, and none is cited here as established. Each has the status of a thing a read source
mentions, which is all it has earned. An entry for any of them would assert a reading that is not
done.

\section{Availability}

The Clay statements, the two papers and the Lean repository are public at the addresses in the
bibliography. The arbitrary-precision arithmetic behind the toolkit chapter is described in the
instruments research paper~\cite{src:Quigg-instruments}.
